\begin{casebox}
\textbf{知识点小案例：CPU 使用率低不等于程序没活干。}
特例是一个写文件管线：业务线程把 block 写入 page cache 后很快返回，但随后 \code{fsync} 阶段 p99 很高。观察上 CPU 使用率可能不高，系统却在等存储设备和内核 writeback。由此推出规律：用户态 runnable、内核等待、page fault、page cache 和设备 completion 是不同边界。工程上不要只用 CPU 利用率判断瓶颈，要把线程状态、系统调用、page fault、dirty bytes、fsync 耗时和设备队列放到同一张阶段表里；不可观测的内核指标要明确降级。
\end{casebox}
